% !TEX program = xelatex
\documentclass[10pt,a4paper]{article}
\newcommand{\coursename}{Математический анализ — подготовка к экзамену}
\newcommand{\rightheader}{Доказательства}
\newcommand{\doctitle}{Доказательства}
% ==== Преамбула экзаменационных материалов (Theory / Proofs / Tasks) ====
% Поля уже, чем в конспектах; после каждого пункта оставлен небольшой
% вертикальный зазор для рукописных пометок.
% Перед этим файлом определяются: \documentclass и макросы
%   \coursename, \rightheader, \doctitle.
\usepackage{fontspec}
\setmainfont{cmunrm.otf}[
  BoldFont       = cmunbx.otf,
  ItalicFont     = cmunti.otf,
  BoldItalicFont = cmunbi.otf]
\usepackage{polyglossia}
\setdefaultlanguage{russian}
\setotherlanguage{english}

\usepackage{amsmath,amssymb,amsthm,mathtools}
\usepackage[a4paper,top=1.8cm,bottom=1.9cm,left=1.7cm,right=1.7cm]{geometry}
\usepackage{enumitem}
% ----- Менее сжатый текст и видимые отступы между абзацами -----
\linespread{1.28}
\setlength{\parskip}{0.55\baselineskip plus 2pt}
\setlist{itemsep=2pt,topsep=3pt,parsep=0pt}
\usepackage{graphicx}
\usepackage{array}
\usepackage{xcolor}
\definecolor{accent}{HTML}{1F6FEB}
\definecolor{teormin}{HTML}{C1121F}
\definecolor{rulecol}{HTML}{D0D7DE}
\usepackage[colorlinks=true,linkcolor=accent,urlcolor=accent,citecolor=accent]{hyperref}
\usepackage{titlesec}
\usepackage{fancyhdr}

% ----- Колонтитулы -----
\pagestyle{fancy}
\fancyhf{}
\fancyhead[L]{\small\itshape \coursename}
\fancyhead[R]{\small\itshape \rightheader}
\fancyfoot[C]{\thepage}
\renewcommand{\headrulewidth}{0.4pt}
\renewcommand{\headrule}{\hbox to\headwidth{\color{rulecol}\leaders\hrule height \headrulewidth\hfill}}

% ----- Заголовки разделов -----
\titleformat{\section}{\large\bfseries\color{accent}}{}{0pt}{}
\titlespacing*{\section}{0pt}{1.2ex}{0.6ex}
\titleformat{\subsection}{\bfseries}{}{0pt}{}
\titlespacing*{\subsection}{0pt}{1.0ex}{0.4ex}

% ----- Операторы и сокращения -----
\DeclareMathOperator{\tg}{tg}
\DeclareMathOperator{\ctg}{ctg}
\DeclareMathOperator{\arctg}{arctg}
\DeclareMathOperator{\arcctg}{arcctg}
\DeclareMathOperator{\sh}{sh}
\DeclareMathOperator{\ch}{ch}
\DeclareMathOperator{\Th}{th}
\DeclareMathOperator{\cth}{cth}
\DeclareMathOperator{\Const}{const}
\DeclareMathOperator{\sign}{sign}
\DeclareMathOperator{\Span}{span}
\newcommand{\R}{\mathbb{R}}
\newcommand{\C}{\mathbb{C}}
\newcommand{\N}{\mathbb{N}}
\newcommand{\Z}{\mathbb{Z}}
\newcommand{\Q}{\mathbb{Q}}
\newcommand{\dx}{\,dx}
\newcommand{\dt}{\,dt}
\newcommand{\du}{\,du}
\newcommand{\eps}{\varepsilon}
\renewcommand{\Re}{\operatorname{Re}}
\renewcommand{\Im}{\operatorname{Im}}
\renewcommand{\le}{\leqslant}
\renewcommand{\ge}{\geqslant}
\newcommand{\scal}[2]{\left( #1,\,#2\right)}
\DeclarePairedDelimiter{\abs}{\lvert}{\rvert}
\DeclarePairedDelimiter{\norm}{\lVert}{\rVert}

% ----- Окружения для доказательств (Proofs) -----
\theoremstyle{definition}
\newtheorem{definition}{Определение}
\newtheorem{example}{Пример}
\newtheorem{remark}{Замечание}
\theoremstyle{plain}
\newtheorem{theorem}{Теорема}
\newtheorem{lemma}{Лемма}
\renewcommand{\proofname}{Доказательство}

% ===== Хелперы для пунктов =====
% Регулируемый зазор для рукописных пометок (пустая строка после пункта):
\newcommand{\writespace}{\vspace{7.5mm}}

% Theory: \faq{номер}{вопрос}{ответ}
% номер теормина помечается красным значком !!! через \tm
\newcommand{\tm}{\textcolor{teormin}{\,!!!}}
\newcommand{\faq}[3]{%
  \par\smallskip\noindent
  \textbf{#1.}\enspace\textit{#2}\par
  \nobreak\vspace{0.3ex}\noindent #3\par
  \writespace}

% Proofs: \prf{номер}{заголовок} ... затем \stmt / \dok блоки
\newcommand{\prf}[2]{%
  \par\writespace\noindent
  {\large\bfseries\color{accent}#1.\enspace #2}\par\smallskip}
\newcommand{\stmt}{\par\noindent\textbf{Формулировка.}\enspace}
\newcommand{\dok}{\par\noindent\textbf{Доказательство.}\enspace}

% Tasks: \task{номер}{условие} ... затем \sol для решения
\newcommand{\task}[2]{%
  \par\writespace\noindent
  \textbf{\textcolor{accent}{Задача #1.}}\enspace #2\par\smallskip}
\newcommand{\sol}{\par\noindent\textbf{Решение.}\enspace}

\setlength{\parindent}{0pt}

\begin{document}
\begin{center}
  {\LARGE\bfseries \doctitle}\\[3pt]
  {\large\itshape \coursename}
\end{center}
\vspace{4pt}

\prf{1}{Интегрирование по частям}
\stmt Если $u(x),v(x)$ непрерывно дифференцируемы, то
$\displaystyle\int u\,dv=uv-\int v\,du$.

\dok Исходный пункт — формула производной произведения:
\[
  (uv)'=u'v+uv'.
\]
Это тождество означает, что $uv$ является первообразной для функции $u'v+uv'$. Поэтому, взяв
неопределённый интеграл от обеих частей, получаем
\[
  uv=\int (u'v+uv')\dx=\int u'v\dx+\int uv'\dx
\]
(интеграл суммы равен сумме интегралов). Перенесём первое слагаемое влево:
\[
  \int uv'\dx=uv-\int u'v\dx .
\]
Осталось заметить, что $v'\dx=dv$ и $u'\dx=du$ — это и есть формула $\int u\,dv=uv-\int v\,du$.
Идея метода: множитель $u$ берут так, чтобы при дифференцировании он \emph{упрощался}, а $dv$ —
так, чтобы он легко интегрировался.

\smallskip
\noindent\textbf{Три типа интегралов.}
\begin{enumerate}[label=\arabic*),leftmargin=*,nosep]
  \item $\int P(x)e^{\alpha x}\dx$, $\int P(x)\sin\beta x\dx$, $\int P(x)\cos\beta x\dx$ — берём
        $u=P(x)$ (многочлен), он понижает степень при каждом дифференцировании.
  \item $\int P(x)\ln x\dx$, $\int P(x)\arctg x\dx$, $\int P(x)\arcsin x\dx$ — берём $u$ равным
        «неудобной» функции ($\ln$, $\arctg$, $\arcsin$): её производная рациональна.
  \item «Циклические» $\int e^{\alpha x}\cos\beta x\dx$ и т.п. — после двукратного применения
        возвращаемся к исходному интегралу и решаем уравнение относительно него.
\end{enumerate}

\prf{2}{Метод замены переменной (подстановки)}
\stmt Пусть $x=\varphi(t)$ непрерывно дифференцируема на промежутке, $f$ имеет первообразную $F$
на множестве значений $\varphi$. Тогда
$\displaystyle\int f(\varphi(t))\varphi'(t)\dt=F(\varphi(t))+C=\int f(x)\dx\Big|_{x=\varphi(t)}$.

\dok Рассмотрим сложную функцию $G(t)=F(\varphi(t))$ и продифференцируем её по правилу цепочки:
\[
  G'(t)=F'(\varphi(t))\cdot\varphi'(t)=f(\varphi(t))\,\varphi'(t),
\]
поскольку $F'=f$ по определению первообразной. Значит $G(t)=F(\varphi(t))$ — первообразная для
функции $f(\varphi(t))\varphi'(t)$, что и означает формулу замены.

\smallskip
\noindent\textbf{Пример (циклический интеграл).} В $\int\dfrac{dx}{\sin x}$ удобна подстановка
$t=\tg\tfrac x2$: тогда $\sin x=\dfrac{2t}{1+t^2}$, $dx=\dfrac{2\,dt}{1+t^2}$, и интеграл
сводится к $\int\dfrac{dt}{t}=\ln\abs{t}+C=\ln\abs{\tg\tfrac x2}+C$.

\prf{3}{Комплексный метод в вещественном анализе}
\stmt $\displaystyle\int e^{ax}\cos bx\dx=\frac{e^{ax}(a\cos bx+b\sin bx)}{a^2+b^2}+C$,\quad
$\displaystyle\int e^{ax}\sin bx\dx=\frac{e^{ax}(a\sin bx-b\cos bx)}{a^2+b^2}+C$.

\dok Рассмотрим комплексную функцию $e^{(a+ib)x}$. По формуле Эйлера
$e^{(a+ib)x}=e^{ax}(\cos bx+i\sin bx)$, поэтому интересующие нас интегралы — это вещественная и
мнимая части одного комплексного интеграла:
\[
  \int e^{(a+ib)x}\dx=\int e^{ax}\cos bx\dx+i\int e^{ax}\sin bx\dx .
\]
Корректность такого разделения: для комплекснозначной функции вещественного аргумента интеграл
определяется покомпонентно, а правило $\bigl(e^{\lambda x}\bigr)'=\lambda e^{\lambda x}$ остаётся
верным и для комплексного $\lambda=a+ib$ (проверяется прямым дифференцированием действительной и
мнимой частей). Поэтому
\[
  \int e^{(a+ib)x}\dx=\frac{e^{(a+ib)x}}{a+ib}+C .
\]
Домножим числитель и знаменатель на сопряжённое $a-ib$, чтобы знаменатель стал вещественным:
\[
  \frac{e^{ax}(\cos bx+i\sin bx)(a-ib)}{(a+ib)(a-ib)}
  =\frac{e^{ax}(\cos bx+i\sin bx)(a-ib)}{a^2+b^2}.
\]
Раскрыв скобки и выделив вещественную часть, получаем формулу для $\int e^{ax}\cos bx\dx$, а мнимую
— для $\int e^{ax}\sin bx\dx$.

\prf{4}{Методы нахождения неопределённых коэффициентов}
\stmt При разложении правильной дроби $\dfrac{P(x)}{Q(x)}$ на простейшие коэффициенты находят
двумя способами; для простого линейного множителя справедлива формула Хевисайда
$A_i=\dfrac{P(a_i)}{Q'(a_i)}$.

\dok После записи разложения с неизвестными коэффициентами обе части умножают на $Q(x)$ и
получают тождество двух многочленов.
\begin{itemize}[leftmargin=*,nosep]
  \item \emph{Метод приравнивания коэффициентов}: приравнивают коэффициенты при одинаковых
        степенях $x$ слева и справа — получается линейная система. Универсален, но громоздок.
  \item \emph{Метод частных значений}: подставляют удобные $x$ (прежде всего корни знаменателя),
        мгновенно обнуляя лишние слагаемые. Быстр для простых корней.
\end{itemize}
\textbf{Обоснование Хевисайда.} Пусть $Q(x)=(x-a_i)\,q(x)$, $q(a_i)\ne0$, и
$\dfrac{P(x)}{Q(x)}=\dfrac{A_i}{x-a_i}+(\text{остальное})$. Умножим на $(x-a_i)$:
\[
  \frac{P(x)}{q(x)}=A_i+(x-a_i)\cdot(\text{остальное}),
\]
и устремим $x\to a_i$ — второе слагаемое исчезает: $A_i=\dfrac{P(a_i)}{q(a_i)}$. Наконец, по
правилу Лейбница $Q'(x)=q(x)+(x-a_i)q'(x)$, откуда $Q'(a_i)=q(a_i)$, и
$A_i=\dfrac{P(a_i)}{Q'(a_i)}$.

\prf{5}{Интегрирование дробей IV типа: рекуррентная формула}
\stmt Для $I_n=\displaystyle\int\frac{dx}{(x^2+a^2)^n}$ ($n\in\N$) справедливо
\[
  I_n=\frac{x}{2a^2(n-1)(x^2+a^2)^{n-1}}+\frac{2n-3}{2a^2(n-1)}\,I_{n-1}.
\]

\dok Отправная точка — тождество, выделяющее $I_{n-1}$:
\[
  I_{n-1}=\int\frac{dx}{(x^2+a^2)^{n-1}}
  =\int\frac{(x^2+a^2)-x^2}{a^2(x^2+a^2)^{n}}\dx
  =\frac{1}{a^2}I_{n-1}-\frac1{a^2}\int\frac{x^2\dx}{(x^2+a^2)^n}.
\]
Оценим последний интеграл интегрированием по частям. Возьмём $u=x$ и
$dv=\dfrac{x\dx}{(x^2+a^2)^n}$; тогда $du=dx$, а
\[
  v=\int\frac{x\dx}{(x^2+a^2)^n}=\frac{1}{2}\int\frac{d(x^2+a^2)}{(x^2+a^2)^n}
   =\frac{(x^2+a^2)^{1-n}}{2(1-n)}=-\frac1{2(n-1)(x^2+a^2)^{n-1}}.
\]
Поэтому
\[
  \int\frac{x^2\dx}{(x^2+a^2)^n}=uv-\int v\,du
  =-\frac{x}{2(n-1)(x^2+a^2)^{n-1}}+\frac1{2(n-1)}I_{n-1}.
\]
Подставляя это в первое равенство и выражая $I_n$, после приведения подобных получаем искомую
рекуррентную формулу. База $I_1=\dfrac1a\arctg\dfrac xa+C$.

\prf{6}{Метод Остроградского}
\stmt Для правильной дроби $\dfrac{P}{Q}$ с кратными корнями знаменателя
\[
  \int\frac{P(x)}{Q(x)}\dx=\frac{P_1(x)}{Q_1(x)}+\int\frac{P_2(x)}{Q_2(x)}\dx,
\]
где $Q_1=\gcd(Q,Q')$, $Q_2=Q/Q_1$, а $P_1,P_2$ — многочлены с неопределёнными коэффициентами
степеней на единицу меньше своих знаменателей.

\dok Идея: $Q_1$ собирает все кратности корней, понизив их на 1, поэтому рациональная (без
логарифмов) часть первообразной имеет знаменатель $Q_1$, а под оставшимся интегралом — только
простые корни (знаменатель $Q_2$). Чтобы найти $P_1,P_2$, продифференцируем основное равенство:
\[
  \frac{P}{Q}=\Bigl(\frac{P_1}{Q_1}\Bigr)'+\frac{P_2}{Q_2}
  =\frac{P_1'Q_1-P_1Q_1'}{Q_1^2}+\frac{P_2}{Q_2}.
\]
Корректность дифференцирования обеспечена тем, что $P_1/Q_1$ дифференцируема всюду, кроме корней
$Q_1$. Приведя правую часть к общему знаменателю $Q=Q_1Q_2$ (используя $Q_1^2\mid Q\cdot Q_1$ через
делимость) и приравняв числители, получаем тождество многочленов; приравнивание коэффициентов при
степенях $x$ даёт линейную систему на коэффициенты $P_1,P_2$.

Метод выгоден, когда $Q$ имеет кратные корни: он сразу выделяет рациональную часть, избавляя от
разложения на простейшие дроби с высокими степенями.

\prf{7}{Интегрирование дробно-линейных иррациональностей}
\stmt Для $R\bigl(x,(\tfrac{ax+b}{cx+d})^{r_1},\dots,(\tfrac{ax+b}{cx+d})^{r_k}\bigr)$ с
рациональными $r_i$ подстановка $t^m=\dfrac{ax+b}{cx+d}$, где $m$ — НОК знаменателей $r_i$,
рационализирует интеграл.

\dok Пусть $r_i=\tfrac{p_i}{q_i}$ и $m=\operatorname{НОК}(q_1,\dots,q_k)$. Положим
$t^m=\dfrac{ax+b}{cx+d}$. Тогда
$\bigl(\tfrac{ax+b}{cx+d}\bigr)^{r_i}=t^{m r_i}=t^{(m/q_i)p_i}$, а $m/q_i\in\Z$, значит показатель
целый — все радикалы исчезли. Выразим $x$: из $t^m(cx+d)=ax+b$ получаем $x(ct^m-a)=b-dt^m$, т.е.
\[
  x=\frac{b-dt^m}{ct^m-a}
\]
— рациональная функция $t$; следовательно $dx$ тоже рациональна. После подстановки подынтегральное
выражение становится рациональной функцией $t$, а такие интегралы берутся в элементарных функциях.

\prf{8}{Квадратичные иррациональности и подстановки Эйлера}
\stmt Каждая из трёх подстановок Эйлера рационализирует $\int R(x,\sqrt{ax^2+bx+c})\dx$.

\dok Разберём \emph{первую} ($a>0$): полагаем $\sqrt{ax^2+bx+c}=\sqrt a\,x+t$. Возведём в квадрат:
\[
  ax^2+bx+c=ax^2+2\sqrt a\,xt+t^2\ \Rightarrow\ bx+c=2\sqrt a\,xt+t^2\ \Rightarrow\
  x=\frac{t^2-c}{b-2\sqrt a\,t}.
\]
Это рациональная функция $t$; тогда и $\sqrt{ax^2+bx+c}=\sqrt a\,x+t$, и $dx$ рациональны по $t$.
Значит $R(x,\sqrt{\dots})\dx$ переходит в рациональную функцию $t$. (\emph{Третья} подстановка при
наличии вещественных корней $x_1,x_2$: $\sqrt{a(x-x_1)(x-x_2)}=t(x-x_1)$ даёт
$a(x-x_2)=t^2(x-x_1)$, откуда $x$ снова рационально выражается через $t$.)

\prf{9}{Интеграл от дифференциального бинома}
\stmt (Чебышёв.) $\int x^m(a+bx^n)^p\dx$ ($m,n,p\in\Q$) берётся в элементарных функциях тогда и
только тогда, когда выполнено одно из условий: (1) $p\in\Z$; (2) $\tfrac{m+1}{n}\in\Z$; (3)
$\tfrac{m+1}{n}+p\in\Z$.

\dok Заменой $z=x^n$ ($x=z^{1/n}$, $dx=\tfrac1n z^{1/n-1}dz$) интеграл приводится к виду
\[
  \frac1n\int z^{\,q}(a+bz)^p\,dz,\qquad q=\frac{m+1}{n}-1 .
\]
\emph{Случай 2:} $\tfrac{m+1}{n}\in\Z\Rightarrow q\in\Z$. Пусть $p=\tfrac us$ ($s$ — знаменатель).
Подстановка $a+bz=w^s$ даёт $z=\tfrac{w^s-a}{b}$ (рационально), а $z^{q}$ — целая степень, тоже
рациональна; всё выражается через $w$ рационально.

\emph{Случай 3:} $\tfrac{m+1}{n}+p\in\Z$. Запишем
$z^{q}(a+bz)^p=z^{q+p}\bigl(\tfrac{a+bz}{z}\bigr)^p$, где $q+p=\tfrac{m+1}{n}-1+p\in\Z$. Подстановка
$\tfrac{a+bz}{z}=az^{-1}+b=w^s$ убирает дробную степень: $z^{-1}=\tfrac{w^s-b}{a}$, $z$ рационально
по $w$, $z^{q+p}$ — целая степень. (\emph{Случай 1:} $p\in\Z$ — подстановка $x=t^N$, $N$ — НОК
знаменателей $m,n$.)

\prf{10}{Определённый интеграл Римана. Необходимое условие интегрируемости}
\stmt $\displaystyle\int_a^b f=\lim_{\lambda(T)\to0}\sigma(f,T,\xi)$. Если $f$ интегрируема, то она
ограничена.

\dok Пусть, напротив, $f$ \emph{не} ограничена на $[a,b]$. При любом разбиении $f$ не ограничена
хотя бы на одном частичном отрезке $[x_{k-1},x_k]$. Зафиксировав все $\xi_i$ при $i\ne k$ и меняя
лишь $\xi_k$ внутри этого отрезка, можно сделать слагаемое $f(\xi_k)\Delta x_k$ сколь угодно
большим по модулю, а значит и всю сумму $\sigma$. Тогда $\sigma$ не имеет конечного предела —
противоречие с интегрируемостью. Следовательно, интегрируемая функция ограничена. (Обратное
неверно: функция Дирихле ограничена, но не интегрируема.)

\prf{11}{Суммы Дарбу при измельчении разбиения}
\stmt При добавлении новых точек нижняя сумма не убывает, верхняя не возрастает; любая нижняя сумма
не превосходит любой верхней.

\dok Добавим одну точку $x'$ в отрезок $[x_{k-1},x_k]$, разбив его на два. Пусть $m_k=\inf$ на всём
отрезке, $m',m''$ — на частях. Поскольку инфимум по меньшему множеству не меньше, $m'\ge m_k$,
$m''\ge m_k$. Вклад в нижнюю сумму меняется с $m_k(x_k-x_{k-1})$ на
$m'(x'-x_{k-1})+m''(x_k-x')\ge m_k(x_k-x_{k-1})$ — не убывает. Аналогично для верхней суммы
($\sup$ по части не больше). Добавление $p$ точек — повторение шага.

Наконец, для двух разбиений $T_1,T_2$ возьмём их объединение $T=T_1\cup T_2$ (измельчение обоих):
$S_*(T_1)\le S_*(T)\le S^*(T)\le S^*(T_2)$.

\prf{12}{Критерий интегрируемости Дарбу}
\stmt $f$ интегрируема $\iff\forall\eps>0\ \exists T:\ S^*(T)-S_*(T)<\eps$ (равносильно
$\underline I=\overline I$).

\dok $(\Leftarrow)$ Из $S^*(T)-S_*(T)<\eps$ при всех $\eps$ следует
$\overline I-\underline I<\eps$, т.е. $\overline I=\underline I=:I$. Для любой выборки
$S_*(T)\le\sigma\le S^*(T)$ и $S_*(T)\le I\le S^*(T)$, поэтому $\abs{\sigma-I}\le S^*(T)-S_*(T)$. По
лемме об измельчении эта разность мала для всех разбиений мельче подходящего $T_0$, значит
$\sigma\to I$.

$(\Rightarrow)$ Если предел $I$ существует, то, выбирая $\xi_i$ так, чтобы $f(\xi_i)$ приближали
$M_i$ (или $m_i$), получаем, что $S^*(T)$ и $S_*(T)$ сами стремятся к $I$, откуда
$S^*(T)-S_*(T)\to0$.

\prf{13}{Интегрируемость непрерывных функций}
\stmt $C[a,b]\subset\mathcal R[a,b]$.

\dok По теореме Кантора $f$ равномерно непрерывна на $[a,b]$: для $\eps>0$ найдётся $\delta>0$, при
котором $\abs{x'-x''}<\delta\Rightarrow\abs{f(x')-f(x'')}<\tfrac{\eps}{b-a}$. Возьмём любое
разбиение с $\lambda(T)<\delta$. На каждом отрезке колебание $\omega_i=M_i-m_i$ достигается как
$f(\xi')-f(\xi'')$ (непрерывность на компакте) и потому $<\tfrac{\eps}{b-a}$. Тогда
\[
  \sum\omega_i\Delta x_i<\frac{\eps}{b-a}\sum\Delta x_i=\eps,
\]
и по критерию через колебания $f$ интегрируема. Теорема Кантора нужна именно для того, чтобы один
$\delta$ годился сразу для всех частичных отрезков.

\prf{14}{Интегрируемость монотонных функций}
\stmt Монотонная на $[a,b]$ функция интегрируема.

\dok Пусть $f$ не убывает. Возьмём разбиение на $n$ равных частей, $\Delta x=\tfrac{b-a}{n}$. В силу
монотонности $m_i=f(x_{i-1})$, $M_i=f(x_i)$, поэтому сумма \emph{телескопическая}:
\[
  \sum\omega_i\Delta x_i=\Delta x\sum_{i=1}^n\bigl(f(x_i)-f(x_{i-1})\bigr)
  =\frac{b-a}{n}\bigl(f(b)-f(a)\bigr)\xrightarrow[n\to\infty]{}0 .
\]
По критерию $f$ интегрируема. Равномерная непрерывность (и теорема Кантора) здесь не нужны:
сокращение происходит за счёт монотонности, поэтому теорема верна даже для функций со счётным
числом разрывов.

\prf{15}{Интегрируемость функций с конечным числом разрывов}
\stmt Ограниченная функция с конечным числом точек разрыва на $[a,b]$ интегрируема.

\dok Пусть $\abs f\le K$ и разрывы в $c_1,\dots,c_p$. Изолируем каждую точку интервалом длины $<\eta$;
их суммарный вклад в $\sum\omega_i\Delta x_i$ не превосходит $2K\cdot p\eta$ (колебание всюду
$\le2K$). Вне этих интервалов $f$ непрерывна на конечном числе замкнутых отрезков, значит
равномерно непрерывна, и там $\sum\omega_i\Delta x_i$ мала (как в №13). Выбрав сперва $\eta$, затем
мелкость разбиения, получаем $\sum\omega_i\Delta x_i<\eps$.

\emph{Ограниченность критична}: для
$1/x$ на $(0,1]$ колебание у нуля бесконечно ($K=\infty$), вклад окрестности не контролируется и
интеграл «не выживает».

\prf{16}{Линейность определённого интеграла}
\stmt $\int_a^b(\alpha f+\beta g)=\alpha\int_a^b f+\beta\int_a^b g$; отсюда — сохранение знака.

\dok Для интегральных сумм $\sigma(\alpha f+\beta g,T,\xi)=\alpha\,\sigma(f,T,\xi)+\beta\,\sigma(g,T,\xi)$
(свойство конечных сумм). Переходя к пределу при $\lambda(T)\to0$ (предел линейной комбинации равен
комбинации пределов), получаем линейность интеграла.

\emph{Сохранение знака:} если $f\ge0$, то все
$\sigma\ge0$, и в пределе $\int_a^b f\ge0$. Применяя это к $g-f\ge0$, получаем монотонность:
$f\le g\Rightarrow\int f\le\int g$.

\prf{17}{Аддитивность интеграла по промежутку}
\stmt $f\in\mathcal R[a,b]\iff f\in\mathcal R[a,c]$ и $f\in\mathcal R[c,b]$ ($a<c<b$), причём
$\int_a^b f=\int_a^c f+\int_c^b f$.

\dok Будем брать разбиения, содержащие точку $c$; такое $T$ распадается на $T'$ (для $[a,c]$) и
$T''$ (для $[c,b]$), и суммы складываются: $S^*(T)=S^*(T')+S^*(T'')$, аналогично нижние. Тогда
\[
  S^*(T)-S_*(T)=\bigl(S^*(T')-S_*(T')\bigr)+\bigl(S^*(T'')-S_*(T'')\bigr),
\]
и обе скобки неотрицательны. Поэтому левая часть мала $\iff$ малы обе скобки — это и есть
эквивалентность интегрируемости на $[a,b]$ и на обеих частях. Переходя к пределу в равенстве сумм,
получаем аддитивность интегралов.

\prf{18}{Оценка модуля интеграла}
\stmt $\abs[\big]{\int_a^b f}\le\int_a^b\abs f$ ($a<b$); при этом $f\in\mathcal R\Rightarrow\abs f\in\mathcal R$.

\dok Сначала интегрируемость $\abs f$: из $\bigl|\,\abs{f(x')}-\abs{f(x'')}\,\bigr|\le\abs{f(x')-f(x'')}$
следует, что колебание $\abs f$ на каждом отрезке не больше колебания $f$, значит
$\sum\omega_i(\abs f)\Delta x_i\le\sum\omega_i(f)\Delta x_i\to0$. Далее из $-\abs f\le f\le\abs f$
интегрированием неравенств: $-\int\abs f\le\int f\le\int\abs f$, т.е. $\abs{\int f}\le\int\abs f$.

\prf{19}{Первая теорема о среднем}
\stmt Для $f\in C[a,b]$ существует $\xi\in[a,b]$: $\int_a^b f=f(\xi)(b-a)$.

\dok Пусть $m=\min f$, $M=\max f$ (достигаются по Вейерштрассу). Из $m\le f\le M$ интегрированием:
$m(b-a)\le\int_a^b f\le M(b-a)$, т.е. среднее $\mu=\tfrac{1}{b-a}\int_a^b f\in[m,M]$. По теореме
Коши о промежуточном значении непрерывная $f$ принимает значение $\mu$: $f(\xi)=\mu$, откуда
$\int_a^b f=f(\xi)(b-a)$. Геометрически — площадь под графиком равна площади прямоугольника высоты
$f(\xi)$.

\prf{20}{Обобщённая теорема о среднем}
\stmt Для $f\in C[a,b]$, $g\in\mathcal R[a,b]$, $g$ знакопостоянной: $\int_a^b fg=f(\xi)\int_a^b g$.

\dok Пусть $g\ge0$. Из $m\le f\le M$ умножением на $g\ge0$: $mg\le fg\le Mg$, интегрируем:
$m\int g\le\int fg\le M\int g$. Если $\int g=0$, то $\int fg=0$ и $\xi$ произвольно. Иначе
$\mu=\tfrac{\int fg}{\int g}\in[m,M]$, и по промежуточному значению $f(\xi)=\mu$.

\emph{Роль знакопостоянства} (контрпример): $f(x)=g(x)=x$ на $[-1,1]$ дают $\int fg=\int x^2=\tfrac23$, но
$\int g=0$, так что равенство $f(\xi)\cdot0=\tfrac23$ невозможно — без знакопостоянства теорема
ложна.

\prf{21}{Непрерывность интеграла с переменным верхним пределом}
\stmt $\Phi(x)=\int_a^x f(t)\dt$ непрерывна на $[a,b]$ при $f\in\mathcal R[a,b]$.

\dok Пусть $\abs f\le K$. Тогда
$\abs{\Phi(x+\Delta x)-\Phi(x)}=\abs[\big]{\int_x^{x+\Delta x}f}\le K\abs{\Delta x}\to0$ при
$\Delta x\to0$. Значит $\Phi$ не просто непрерывна, а удовлетворяет условию Липшица с константой
$K$.

\prf{22}{Теорема Барроу}
\stmt Если $f$ непрерывна в точке $x$, то $\Phi'(x)=f(x)$.

\dok Рассмотрим
\[
  \frac{\Phi(x+\Delta x)-\Phi(x)}{\Delta x}-f(x)
  =\frac1{\Delta x}\int_x^{x+\Delta x}\bigl(f(t)-f(x)\bigr)\dt .
\]
Оценим по модулю: правая часть $\le\sup_{\abs{t-x}\le\abs{\Delta x}}\abs{f(t)-f(x)}$. В силу
непрерывности $f$ \emph{в точке} $x$ этот супремум стремится к нулю при $\Delta x\to0$.
Следовательно $\Phi'(x)=f(x)$. Именно на этом шаге (и только в исследуемой точке) нужна
непрерывность подынтегральной функции.

\prf{23}{Приложение теоремы Барроу к пределам (правило Лопиталя)}
\stmt Неопределённости $\tfrac00$, $\tfrac\infty\infty$ для отношений интегралов с переменными
пределами раскрываются дифференцированием по Барроу.

\dok Пусть $F(x)=\int_{a}^{\varphi(x)}f$, $G(x)=\int_{a}^{\psi(x)}h$, и обе $\to0$ при $x\to x_0$.
По теореме Барроу и правилу цепочки $F'(x)=f(\varphi(x))\varphi'(x)$,
$G'(x)=h(\psi(x))\psi'(x)$. Если $\varphi,\psi$ дифференцируемы, $G'\ne0$ в проколотой окрестности
$x_0$ и существует $\lim\tfrac{F'}{G'}$, то по правилу Лопиталя (теорема Коши о среднем)
$\lim\tfrac FG=\lim\tfrac{F'}{G'}$.

\prf{24}{Теорема Ньютона–Лейбница}
\stmt Для $f\in C[a,b]$ и любой первообразной $F$: $\int_a^b f=F(b)-F(a)$.

\dok По теореме Барроу $\Phi(x)=\int_a^x f$ — первообразная для $f$. Любая другая первообразная
отличается на константу: $F(x)=\Phi(x)+C$. Тогда
\[
  F(b)-F(a)=\Phi(b)-\Phi(a)=\int_a^b f-\int_a^a f=\int_a^b f .
\]
Значение теоремы: она заменяет вычисление пределов сложных интегральных сумм нахождением одной
первообразной.

\prf{25}{Площадь криволинейного сектора в полярных координатах}
\stmt Для $\rho\in C[\alpha,\beta]$ площадь $S=\tfrac12\int_\alpha^\beta\rho^2(\varphi)\,d\varphi$.

\dok Разобьём $[\alpha,\beta]$. На $[\varphi_{i-1},\varphi_i]$ заключим часть сектора между двумя
круговыми секторами радиусов $m_i=\min\rho$ и $M_i=\max\rho$:
$\tfrac12 m_i^2\Delta\varphi_i\le\Delta S_i\le\tfrac12 M_i^2\Delta\varphi_i$. Суммы по $i$ —
это нижняя и верхняя суммы Дарбу для функции $\tfrac12\rho^2$. Так как $\rho^2$ непрерывна (а значит
интегрируема), обе суммы сходятся к $\tfrac12\int_\alpha^\beta\rho^2\,d\varphi$, к нему же зажата и
площадь $S$.

\prf{26}{Длина дуги гладкой кривой}
\stmt Для $y=y(x)$, $y'\in C[a,b]$: $L=\int_a^b\sqrt{1+y'^2}\dx$.

\dok Длина вписанной ломаной по разбиению равна
\[
  \sum_i\sqrt{\Delta x_i^2+\Delta y_i^2}=\sum_i\sqrt{1+\Bigl(\tfrac{\Delta y_i}{\Delta x_i}\Bigr)^2}\,
  \Delta x_i .
\]
По теореме Лагранжа о конечном приращении $\tfrac{\Delta y_i}{\Delta x_i}=y'(\xi_i)$ для некоторого
$\xi_i\in[x_{i-1},x_i]$. Получаем интегральную сумму
$\sum\sqrt{1+y'(\xi_i)^2}\,\Delta x_i$ для непрерывной (а значит интегрируемой) функции
$\sqrt{1+y'^2}$. Предел при $\lambda(T)\to0$ равен $\int_a^b\sqrt{1+y'^2}\dx$.

\prf{27}{Объёмы тел вращения: диски и оболочки}
\stmt Для $y=f(x)\ge0$, $f\in C[a,b]$: вокруг $Ox$ — $V=\pi\int_a^b f^2\dx$; вокруг $Oy$ ($a\ge0$) —
$V=2\pi\int_a^b x f(x)\dx$.

\dok (a) Сечение тела плоскостью $x=\Const$ — круг радиуса $f(x)$, его площадь $\pi f^2(x)$. Объём
слоя толщины $\Delta x$ приближается элементарным диском $\pi f(\xi_i)^2\Delta x_i$; сумма — это
интегральная сумма для $\pi f^2$, и в пределе $V=\pi\int_a^b f^2\dx$.

(b) Кольцевой слой $[x_{i-1},x_i]$ при вращении вокруг $Oy$ даёт цилиндрическую оболочку радиуса
$x$, высоты $f(x)$, толщины $\Delta x$: её объём $\approx2\pi x_i f(\xi_i)\Delta x_i$. Сумма —
интегральная для $2\pi x f(x)$, предел $V=2\pi\int_a^b x f(x)\dx$.

\prf{28}{Признаки Дирихле и Абеля для несобственных интегралов}
\stmt (Дирихле.) $\int_a^{+\infty}f g\dx$ сходится, если $F(A)=\int_a^A f$ ограничена, а $g$
монотонно стремится к нулю.

\dok По второй теореме о среднем (форма Бонне) на $[A_1,A_2]$:
$\int_{A_1}^{A_2}fg=g(A_1)\int_{A_1}^{\eta}f+g(A_2)\int_{\eta}^{A_2}f$ для некоторого $\eta$.
Поскольку $\abs{\int_{p}^{q}f}=\abs{F(q)-F(p)}\le2M$ (ограниченность $F$),
\[
  \abs[\Big]{\int_{A_1}^{A_2}fg}\le2M\bigl(\abs{g(A_1)}+\abs{g(A_2)}\bigr)\xrightarrow[A_1,A_2\to\infty]{}0,
\]
т.е. выполнен критерий Коши — интеграл сходится.

\emph{Пример:} $\int_0^{+\infty}\tfrac{\sin x}{x}\dx$
с $f=\sin x$ ($\abs{\int_0^A\sin}\le2$) и $g=\tfrac1x\downarrow0$ сходится. Абсолютной сходимости
нет: $\int_0^{+\infty}\tfrac{\abs{\sin x}}{x}\ge\sum_{k}\tfrac1{(k+1)\pi}\int_{k\pi}^{(k+1)\pi}\abs{\sin x}\dx
=\sum_k\tfrac{2}{(k+1)\pi}=\infty$.

\prf{29}{Определение ряда и необходимый признак сходимости}
\stmt Ряд сходится, если существует конечный $\lim S_N=S$. Необходимый признак: $a_n\to0$.

\dok Если $S_N\to S$, то $a_n=S_n-S_{n-1}\to S-S=0$. Обратное неверно: для гармонического ряда
$\tfrac1n\to0$, однако
$S_{2n}-S_n=\sum_{k=n+1}^{2n}\tfrac1k\ge n\cdot\tfrac1{2n}=\tfrac12$, значит $\{S_n\}$ не
фундаментальна и ряд расходится.

\prf{30}{Критерий Коши для числового ряда}
\stmt Ряд сходится $\iff\forall\eps>0\ \exists N\ \forall n>N\ \forall p:\
\abs{a_{n+1}+\dots+a_{n+p}}<\eps$.

\dok Сходимость ряда означает существование предела $\{S_N\}$, что по критерию Коши для
последовательностей (полнота $\R$) равносильно фундаментальности:
$\abs{S_{n+p}-S_n}=\abs{a_{n+1}+\dots+a_{n+p}}<\eps$. Для гармонического ряда возьмём $p=n$:
$\abs{S_{2n}-S_n}\ge\tfrac12$ — условие нарушено, ряд расходится.

\prf{31}{Линейные свойства рядов и конечные изменения}
\stmt Линейность сходящихся рядов; отбрасывание конечного числа членов не влияет на факт
сходимости.

\dok Если $S_n^a\to A$, $S_n^b\to B$, то $\alpha S_n^a+\beta S_n^b\to\alpha A+\beta B$, т.е.
$\sum(\alpha a_n+\beta b_n)=\alpha A+\beta B$. Отбросив первые $m$ членов, для $n>m$ изменяем $S_n$
на постоянную величину (сумму отброшенных); предел существует или нет одновременно с исходным,
поэтому факт сходимости инвариантен, а сумма сдвигается.

\prf{32}{Ряды с неотрицательными членами и признаки сравнения}
\stmt $a_n\ge0$: ряд сходится $\iff$ частичные суммы ограничены сверху. Признаки сравнения в форме
неравенств и в предельной форме.

\dok При $a_n\ge0$ суммы $S_n$ не убывают, значит сходятся $\iff$ ограничены сверху (теорема о
монотонной ограниченной последовательности). \emph{Неравенства:} $0\le a_n\le b_n\Rightarrow
S_n^a\le S_n^b$; ограниченность $S^b$ влечёт ограниченность $S^a$.

\emph{Предельная форма:} если
$\tfrac{a_n}{b_n}\to L\in(0,\infty)$, то при больших $n$ $c_1 b_n\le a_n\le c_2 b_n$, и ряды ведут
себя одинаково. Пример: $\sum\tfrac{\ln n}{n^2+\sin n}\sim\tfrac{\ln n}{n^2}\le\tfrac1{n^{3/2}}$ при
больших $n$ — сходится.

\prf{33}{Признак Д'Аламбера}
\stmt $\lim\tfrac{a_{n+1}}{a_n}=q$: $q<1$ сходится, $q>1$ расходится, $q=1$ — нет ответа.

\dok Пусть $q<1$; возьмём $q<r<1$. Существует $N$, при котором $\tfrac{a_{n+1}}{a_n}<r$ для $n\ge N$,
откуда $a_{N+k}<a_N r^k$. Ряд мажорируется сходящейся прогрессией $a_N\sum r^k$ — сходится. При
$q>1$ члены растут, $a_n\not\to0$ — расходимость. При $q=1$ признак бессилен: $\sum\tfrac1n$
расходится, $\sum\tfrac1{n^2}$ сходится, у обоих $q=1$.

\prf{34}{Радикальный признак Коши}
\stmt $\lim\sqrt[n]{a_n}=q$: $q<1$ сходится, $q>1$ расходится. Коши сильнее Д'Аламбера.

\dok $q<1$: при $q<r<1$ имеем $\sqrt[n]{a_n}<r$, т.е. $a_n<r^n$ для больших $n$ — мажоранта
$\sum r^n$. $q>1$: $a_n>1$ бесконечно часто, $a_n\not\to0$.

\emph{Сила:} если существует
$\lim\tfrac{a_{n+1}}{a_n}=q$, то существует и $\lim\sqrt[n]{a_n}=q$ (предел среднего
геометрического / теорема Штольца), но обратное неверно — есть ряды, где работает Коши, а
Д'Аламбер не имеет предела.

\prf{35}{Абсолютная сходимость и её свойства}
\stmt Абсолютная сходимость влечёт сходимость.

\dok Пусть $\sum\abs{a_n}$ сходится. По критерию Коши $\abs{a_{n+1}}+\dots+\abs{a_{n+p}}<\eps$ при
$n>N$. По неравенству треугольника $\abs{a_{n+1}+\dots+a_{n+p}}\le\abs{a_{n+1}}+\dots+\abs{a_{n+p}}<\eps$,
т.е. выполнен критерий Коши и для $\sum a_n$ — ряд сходится. (Без доказательства: теорема Римана —
условно сходящийся ряд перестановкой можно привести к любой сумме, чего не бывает у абсолютно
сходящихся.)

\prf{36}{Условная сходимость и знакочередующийся гармонический ряд}
\stmt $\sum\tfrac{(-1)^{n-1}}{n}$ сходится условно.

\dok Ряд из модулей — гармонический $\sum\tfrac1n$ — расходится (№30), значит абсолютной сходимости
нет. Сам ряд знакочередующийся с $c_n=\tfrac1n\downarrow0$, поэтому по признаку Лейбница (№37)
сходится. Сходимость без абсолютной — это и есть условная сходимость.

\prf{37}{Признак Лейбница и оценка остатка}
\stmt Если $c_n>0$ монотонно убывают к нулю, то $\sum(-1)^{n-1}c_n$ сходится и $\abs{R_N}\le c_{N+1}$.

\dok Чётные суммы $S_{2k}=(c_1-c_2)+\dots+(c_{2k-1}-c_{2k})$ не убывают (скобки $\ge0$). С другой
стороны $S_{2k}=c_1-(c_2-c_3)-\dots-c_{2k}<c_1$ — ограничены сверху, значит $S_{2k}\to S$. Так как
$S_{2k+1}=S_{2k}+c_{2k+1}\to S+0=S$, то и $S_n\to S$. Геометрически частичные суммы колеблются
вокруг $S$ с убывающей амплитудой.

Остаток $R_N$ — снова лейбницевский ряд с первым членом
$c_{N+1}$, поэтому $\abs{R_N}\le c_{N+1}$ (и знак $R_N$ совпадает со знаком первого отброшенного
члена).

\prf{38}{Преобразование Абеля (суммирование по частям)}
\stmt $\displaystyle\sum_{k=1}^n a_k b_k=a_n B_n-\sum_{k=1}^{n-1}(a_{k+1}-a_k)B_k$, где
$B_k=\sum_{j=1}^k b_j$, $B_0=0$.

\dok Так как $b_k=B_k-B_{k-1}$,
\[
  \sum_{k=1}^n a_k b_k=\sum_{k=1}^n a_k(B_k-B_{k-1})
  =\sum_{k=1}^n a_k B_k-\sum_{k=1}^n a_k B_{k-1}.
\]
Во второй сумме сдвинем индекс ($k\mapsto k+1$, учитывая $B_0=0$):
$\sum_{k=1}^n a_k B_{k-1}=\sum_{k=1}^{n-1}a_{k+1}B_k$. Вычитая, получаем
$\sum a_k b_k=a_n B_n+\sum_{k=1}^{n-1}(a_k-a_{k+1})B_k$ — тождество Абеля (дискретный аналог
интегрирования по частям).

\prf{39}{Признак Дирихле сходимости рядов}
\stmt $\sum a_n b_n$ сходится, если частичные суммы $B_n=\sum b_k$ ограничены ($\abs{B_n}\le M$), а
$a_n\downarrow0$ монотонно.

\dok По преобразованию Абеля для «хвоста» и лемме об оценке частичных сумм
$\abs[\big]{\sum_{k=n+1}^{n+p}a_k b_k}\le 2M\,a_{n+1}$. Так как $a_{n+1}\to0$, правая часть $<\eps$
при больших $n$ — выполнен критерий Коши, ряд сходится.

\emph{Пример} $\sum\tfrac{\sin n}{n^p}$:
суммы $\sum_{k\le n}\sin k$ ограничены (замкнутая формула, $\le\tfrac1{\abs{\sin(1/2)}}$),
$a_n=\tfrac1{n^p}\downarrow0$ при $p>0$ — ряд сходится. При $0<p\le1$ сходимость лишь условная
($\tfrac{\abs{\sin n}}{n^p}$ в среднем $\sim\tfrac{2/\pi}{n^p}$ даёт расходящийся ряд из модулей).

\prf{40}{Признак Абеля сходимости рядов}
\stmt $\sum a_n b_n$ сходится, если $\sum b_n$ сходится, а $\{a_n\}$ монотонна и ограничена.

\dok Монотонная ограниченная $\{a_n\}$ имеет предел $a$. Запишем $a_n=a+\alpha_n$, где
$\alpha_n=a_n-a$ монотонно стремится к нулю. Тогда $\sum a_n b_n=a\sum b_n+\sum\alpha_n b_n$. Первый
ряд сходится (по условию). Второй — по признаку Дирихле: частичные суммы $\sum b_k$ ограничены (как
у сходящегося ряда), а $\alpha_n\downarrow0$. Сходимость базового $\sum b_n$ как раз и позволяет не
требовать $a_n\to0$.

\prf{41}{Поточечная и равномерная сходимость функций}
\stmt Определения поточечной и равномерной сходимости; геометрический смысл.

\dok \emph{Поточечно}: $\forall x\ \forall\eps\ \exists N(\eps,x)\ \forall n>N:\
\abs{f_n(x)-f(x)}<\eps$ — для каждого $x$ свой номер. \emph{Равномерно}: тот же $N(\eps)$ годится
для всех $x$ сразу ($\exists N$ \emph{до} $\forall x$).

Геометрически равномерная сходимость
означает, что начиная с номера $N$ весь график $f_n$ лежит в «трубке»
$\{(x,y):\abs{y-f(x)}<\eps\}$ вокруг предельной функции. При поточечной сходимости момент входа в
трубку зависит от $x$, и единого $N$ может не быть.

\prf{42}{Супремум-критерий равномерной сходимости}
\stmt $f_n\rightrightarrows f\iff\rho_n=\sup_X\abs{f_n-f}\to0$.

\dok $(\Leftarrow)$ Если $\rho_n\to0$, то $\abs{f_n(x)-f(x)}\le\rho_n<\eps$ для всех $x$ при $n>N$ —
равномерно. $(\Rightarrow)$ Если равномерно, то для $\eps$ найдётся $N$: $\abs{f_n(x)-f(x)}<\eps$
всюду, откуда $\rho_n\le\eps$, т.е. $\rho_n\to0$.

\emph{Бегущий горб:} $f_n(x)=x^n-x^{2n}$ на
$[0,1]$ поточечно $\to0$, но максимум при $x^n=\tfrac12$ равен $\tfrac14$, значит
$\rho_n=\tfrac14\not\to0$ — равномерной сходимости нет.

\prf{43}{Критерий Коши для функциональных последовательностей}
\stmt $\{f_n\}$ сходится равномерно на $X\iff\forall\eps\ \exists N\ \forall n>N\ \forall p\
\forall x:\ \abs{f_{n+p}(x)-f_n(x)}<\eps$.

\dok $(\Rightarrow)$ Если $f_n\rightrightarrows f$, то
$\abs{f_{n+p}(x)-f_n(x)}\le\abs{f_{n+p}-f}+\abs{f-f_n}<\eps$ (неравенство треугольника).

$(\Leftarrow)$ При каждом фиксированном $x$ последовательность $\{f_n(x)\}$ фундаментальна, значит
(полнота $\R$) сходится к некоторому $f(x)$. Перейдём в неравенстве $\abs{f_{n+p}(x)-f_n(x)}<\eps$ к
пределу $p\to\infty$: $\abs{f(x)-f_n(x)}\le\eps$ для всех $x$ при $n>N$ — сходимость равномерна.

\prf{44}{Мажорантный признак Вейерштрасса}
\stmt Если $\abs{u_n(x)}\le a_n$ для всех $x$ и $\sum a_n$ сходится, то $\sum u_n$ сходится
равномерно (и абсолютно).

\dok По критерию Коши для $\sum a_n$: $\sum_{k=n+1}^{n+p}a_k<\eps$ при $n>N$. Тогда
$\abs[\big]{\sum_{k=n+1}^{n+p}u_k(x)}\le\sum_{k=n+1}^{n+p}\abs{u_k(x)}\le\sum_{k=n+1}^{n+p}a_k<\eps$
равномерно по $x$ — выполнен критерий Коши для функционального ряда.

Признак достаточный, но не
необходимый: существуют равномерно сходящиеся ряды (например, с членами $\tfrac{(-1)^n}{n}$,
не зависящими от $x$), которые нельзя мажорировать сходящимся числовым рядом, ибо $\sum\tfrac1n=\infty$.

\prf{45}{Непрерывность суммы ряда (метод трёх $\eps$)}
\stmt Если $u_n$ непрерывны в окрестности $x_0$ и $\sum u_n\rightrightarrows S$, то $S$ непрерывна в
$x_0$.

\dok Запишем $S=S_N+R_N$ и оценим
\[
  \abs{S(x)-S(x_0)}\le\abs{S_N(x)-S_N(x_0)}+\abs{R_N(x)}+\abs{R_N(x_0)} .
\]
По равномерной сходимости выберем $N$ так, что $\abs{R_N}<\tfrac\eps3$ всюду (сразу для второго и
третьего слагаемых). При этом фиксированном $N$ сумма $S_N$ — конечная сумма непрерывных функций,
значит непрерывна: $\exists\delta>0$, при $\abs{x-x_0}<\delta$ имеем
$\abs{S_N(x)-S_N(x_0)}<\tfrac\eps3$. Итого $\abs{S(x)-S(x_0)}<\eps$.

\prf{46}{Почленное интегрирование и дифференцирование рядов}
\stmt Если $u_n\in C[a,b]$ и $\sum u_n\rightrightarrows S$, то $\int_a^b S=\sum\int_a^b u_n$.

\dok $S$ непрерывна (№45), значит интегрируема. Оценим разность интегралов:
\[
  \abs[\Big]{\int_a^b S-\int_a^b S_n}=\abs[\Big]{\int_a^b R_n}\le\int_a^b\abs{R_n}
  \le(b-a)\sup_{[a,b]}\abs{R_n}\to0 .
\]
Значит $\int_a^b S_n\to\int_a^b S$, т.е. $\int_a^b S=\lim\sum_{k\le n}\int_a^b u_k=\sum\int_a^b u_n$.
(Теорема о почленном дифференцировании выводится отсюда применением к равномерно сходящемуся ряду
производных.)

\prf{47}{Степенной ряд и первая теорема Абеля}
\stmt Если $\sum a_n x^n$ сходится в $x_0\ne0$, то он абсолютно сходится при $\abs x<\abs{x_0}$;
отсюда — существование радиуса сходимости.

\dok Из сходимости следует $a_n x_0^n\to0$, значит $\abs{a_n x_0^n}\le M$. При $\abs x<\abs{x_0}$ и
$q=\abs{x/x_0}<1$:
\[
  \abs{a_n x^n}=\abs{a_n x_0^n}\cdot q^n\le M q^n,
\]
и ряд мажорируется сходящейся прогрессией — абсолютная сходимость. Поэтому множество точек
сходимости есть интервал; его «радиус» $R=\sup\{\abs x:\text{сходится}\}$ симметричен относительно
центра: на $(-R,R)$ ряд сходится, вне — расходится.

\prf{48}{Теорема Коши–Адамара}
\stmt $\dfrac1R=\varlimsup\limits_{n\to\infty}\sqrt[n]{\abs{a_n}}$.

\dok Применим радикальный признак к ряду модулей $\sum\abs{a_n x^n}$:
\[
  \varlimsup\sqrt[n]{\abs{a_n x^n}}=\abs x\cdot\varlimsup\sqrt[n]{\abs{a_n}}=\frac{\abs x}{R}.
\]
Если $\abs x<R$, предел $<1$ — абсолютная сходимость; если $\abs x>R$, то $\sqrt[n]{\abs{a_n x^n}}>1$
бесконечно часто, $a_n x^n\not\to0$ — расходимость (нарушен необходимый признак). Случаи $R=0$ и
$R=+\infty$ отвечают $\varlimsup=+\infty$ и $0$.

\prf{49}{Свойства степенного ряда внутри интервала сходимости}
\stmt На любом $[-r,r]$, $0<r<R$, ряд сходится абсолютно и равномерно; отсюда непрерывность суммы,
почленное интегрирование и (бесконечное) дифференцирование.

\dok При $\abs x\le r$: $\abs{a_n x^n}\le\abs{a_n}r^n$, а ряд $\sum\abs{a_n}r^n$ сходится ($r<R$). По
признаку Вейерштрасса сходимость на $[-r,r]$ равномерна. Из равномерной сходимости получаем
непрерывность суммы и право почленно интегрировать по любому внутреннему отрезку.

Ряд производных
$\sum n a_n x^{n-1}$ имеет тот же радиус $R$ (так как $\sqrt[n]{n}\to1$) и потому тоже равномерно
сходится на $[-r,r]$ — значит ряд можно дифференцировать почленно, причём сколько угодно раз.

\prf{50}{Вещественные аналитические функции и пример Коши}
\stmt $C^\omega\subsetneq C^\infty$: функция $f(x)=e^{-1/x^2}$ ($f(0)=0$) гладкая, но не
аналитична в нуле.

\dok При $x\ne0$ каждая производная имеет вид $f^{(n)}(x)=P_n(1/x)\,e^{-1/x^2}$ (многочлен от
$1/x$). В нуле по индукции $f^{(n)}(0)=\lim_{x\to0}P_n(1/x)e^{-1/x^2}=0$, поскольку $e^{-1/x^2}$
убывает быстрее любой степени $1/x$. Значит все $f^{(n)}(0)=0$ и ряд Тейлора в нуле тождественно
равен нулю — он не совпадает с $f(x)\ne0$ ни в одной проколотой окрестности.

Причина «ложной гладкости»: гладкость не ограничивает рост производных, и ряд Тейлора не
восстанавливает функцию.

\prf{51}{Коэффициенты Тейлора и единственность разложения}
\stmt Если $f(x)=\sum a_n(x-x_0)^n$ на $(x_0-R,x_0+R)$, то $a_k=\dfrac{f^{(k)}(x_0)}{k!}$, и
представление единственно.

\dok Почленно дифференцируем $k$ раз (внутри интервала это законно, №49):
\[
  f^{(k)}(x)=\sum_{n\ge k}a_n\,n(n-1)\cdots(n-k+1)(x-x_0)^{n-k}.
\]
Подставив $x=x_0$, обнуляем все слагаемые с $n>k$; остаётся член $n=k$: $f^{(k)}(x_0)=a_k\,k!$,
откуда $a_k=f^{(k)}(x_0)/k!$. Так как коэффициенты однозначно выражаются через производные суммы,
два степенных ряда с одинаковой суммой имеют равные коэффициенты — представление единственно.

\prf{52}{Критерий аналитичности (ограничение роста производных)}
\stmt Если в окрестности $x_0$ выполнено $\abs{f^{(n)}(x)}\le C\,M^n$ для всех $n$, то $f$
аналитична.

\dok Остаточный член в форме Лагранжа: при некотором $\xi$ между $x_0$ и $x$
\[
  \abs{R_n(x)}=\abs[\Big]{\frac{f^{(n+1)}(\xi)}{(n+1)!}(x-x_0)^{n+1}}
  \le C\,\frac{M^{n+1}\abs{x-x_0}^{n+1}}{(n+1)!}\xrightarrow[n\to\infty]{}0,
\]
поскольку факториал в знаменателе растёт быстрее любой геометрической прогрессии. Значит ряд Тейлора
сходится к $f$ — функция аналитична.

\prf{53}{Евклидовы пространства и неравенство Коши–Буняковского}
\stmt $\abs{(f,g)}\le\norm f\cdot\norm g$ в предгильбертовом $L_2[a,b]$.

\dok При любом $\lambda\in\R$ квадратный относительно $\lambda$ трёхчлен
$\norm{\lambda f+g}^2=\lambda^2\norm f^2+2\lambda(f,g)+\norm g^2\ge0$ неотрицателен, поэтому его
дискриминант $\le0$: $4(f,g)^2-4\norm f^2\norm g^2\le0$, т.е. $\abs{(f,g)}\le\norm f\norm g$.
\emph{Фактор-пространство:} $\norm f=0$ не влечёт $f\equiv0$ ($f$ может быть отлична от нуля на
множестве меры нуль); отождествляя функции, совпадающие почти всюду, делаем $\norm{\cdot}$
корректной нормой.

\prf{54}{Экстремальное свойство коэффициентов Фурье. Неравенство Бесселя}
\stmt Среди сумм $\sum_{n\le N}\alpha_n\varphi_n$ среднеквадратичное отклонение минимально при
$\alpha_n=c_n$; отсюда $\sum c_n^2\norm{\varphi_n}^2\le\norm f^2$.

\dok В силу ортогональности
\[
  \delta_N=\norm[\Big]{f-\sum_{n\le N}\alpha_n\varphi_n}^2
  =\norm f^2-2\sum\alpha_n(f,\varphi_n)+\sum\alpha_n^2\norm{\varphi_n}^2 .
\]
Подставив $(f,\varphi_n)=c_n\norm{\varphi_n}^2$ и выделив полный квадрат,
\[
  \delta_N=\norm f^2-\sum_{n\le N}c_n^2\norm{\varphi_n}^2+\sum_{n\le N}(\alpha_n-c_n)^2\norm{\varphi_n}^2 .
\]
Последняя сумма $\ge0$ и минимальна (равна $0$) при $\alpha_n=c_n$. Тогда $\delta_N\ge0$ даёт
$\sum_{n\le N}c_n^2\norm{\varphi_n}^2\le\norm f^2$; переход $N\to\infty$ — неравенство Бесселя.

\prf{55}{Теорема Стеклова (полнота $\iff$ замкнутость)}
\stmt Ортогональная система замкнута (равенство Парсеваля для всех $f$) $\iff$ полна.

\dok $(\Rightarrow)$ Пусть система замкнута и $(f,\varphi_n)=0$ для всех $n$, т.е. все $c_n=0$. По
Парсевалю $\norm f^2=\sum c_n^2\norm{\varphi_n}^2=0$, значит $f\equiv0$ — система полна.

$(\Leftarrow)$ Пусть система полна. Остаток $R_N=f-\sum_{n\le N}c_n\varphi_n$; по Бесселю
$\sum c_n^2\norm{\varphi_n}^2$ сходится, поэтому $\{S_N\}$ фундаментальна и (полнота $L_2$) сходится
к некоторой $g$. Разность $f-g$ ортогональна всем $\varphi_n$, а в силу полноты $f-g\equiv0$, т.е.
$S_N\to f$. Тогда $\norm{R_N}^2=\norm f^2-\sum_{n\le N}c_n^2\norm{\varphi_n}^2\to0$ — равенство
Парсеваля.

\prf{56}{Лемма Римана–Лебега}
\stmt Для интегрируемой $f$: $\int_a^b f(x)\sin(\lambda x)\dx\to0$ при $\lambda\to\infty$
(аналогично с $\cos$).

\dok \emph{Шаг 1 — ступенчатая функция.} Если $f$ постоянна на $[c,d]$, то
$\int_c^d\sin\lambda x\dx=\tfrac{\cos\lambda c-\cos\lambda d}{\lambda}\to0$; конечная сумма таких
интегралов тоже стремится к нулю.

\emph{Шаг 2 — общий случай.} Интегрируемую $f$ приблизим
ступенчатой $g$ так, что $\int_a^b\abs{f-g}<\eps$. Тогда
\[
  \abs[\Big]{\int_a^b f\sin\lambda x}\le\int_a^b\abs{f-g}+\abs[\Big]{\int_a^b g\sin\lambda x}
  <\eps+\bigl(\to0\bigr),
\]
поэтому верхний предел не превосходит $\eps$; в силу произвольности $\eps$ интеграл стремится к
нулю.

\prf{57}{Принцип локализации Римана}
\stmt Сходимость ряда Фурье в точке $x_0$ и его сумма зависят лишь от поведения $f$ в сколь угодно
малой окрестности $x_0$.

\dok Используя интегральное представление, $S_N(x_0)-S=\tfrac1\pi\int_{-\pi}^{\pi}(f(x_0+t)-S)D_N(t)\dt$,
где $D_N(t)=\tfrac{\sin((N+\frac12)t)}{2\sin(t/2)}$. Вне окрестности $\abs t<\delta$ знаменатель
$2\sin\tfrac t2$ отделён от нуля, поэтому функция $\tfrac{f(x_0+t)-S}{2\sin(t/2)}$ интегрируема, и по
лемме Римана–Лебега $\int_{\delta\le\abs t\le\pi}(\dots)\sin((N+\tfrac12)t)\dt\to0$. Значит вклад в
предел даёт только интеграл по $\abs t<\delta$ — поведение вне окрестности несущественно.

\prf{58}{Признак Дини и условие Липшица}
\stmt Если $\int_0^\pi\tfrac{\abs{f(x_0+t)+f(x_0-t)-2S}}{t}\dt<\infty$, ряд Фурье сходится в $x_0$ к
$S$. Условие Липшица $\Rightarrow$ Дини.

\dok В силу чётности ядра
$S_N(x_0)-S=\tfrac1\pi\int_0^\pi\bigl(f(x_0+t)+f(x_0-t)-2S\bigr)D_N(t)\dt$. Введём
$g(t)=\tfrac{f(x_0+t)+f(x_0-t)-2S}{2\sin(t/2)}$. Поскольку $\sin\tfrac t2\sim\tfrac t2$, условие Дини
гарантирует интегрируемость $g$, и по лемме Римана–Лебега
$\int_0^\pi g(t)\sin\bigl((N+\tfrac12)t\bigr)\dt\to0$, т.е. $S_N(x_0)\to S$.

\emph{Липшиц:} если
$\abs{f(x_0\pm t)-f(x_0)}\le Lt^\alpha$, берём $S=f(x_0)$; тогда числитель $O(t^\alpha)$ и
$\tfrac{\abs{\dots}}{t}\le2L\,t^{\alpha-1}$ интегрируема в нуле ($\alpha>0$). Дифференцируемая
функция удовлетворяет Липшицу с $\alpha=1$, значит её ряд Фурье сходится.

\prf{59}{Ряды Фурье в точках разрыва первого рода}
\stmt Для кусочно-гладкой $f$ ряд Фурье в точке разрыва I рода сходится к
$\tfrac{f(x_0-0)+f(x_0+0)}{2}$.

\dok Берём $S=\tfrac{f(x_0-0)+f(x_0+0)}{2}$. Для кусочно-гладкой $f$ справа
$f(x_0+t)-f(x_0+0)=O(t)$, слева $f(x_0-t)-f(x_0-0)=O(t)$, поэтому $f(x_0+t)+f(x_0-t)-2S=O(t)$ и
подынтегральное выражение в условии Дини ограничено — условие выполнено, $S_N\to S$.

\emph{Пример} $f(x)=\sign x$ на $[-\pi,\pi]$: $a_n=0$,
$b_n=\tfrac1\pi\int_{-\pi}^\pi\sign x\sin nx\dx=\tfrac2\pi\int_0^\pi\sin nx\dx=\tfrac{2(1-(-1)^n)}{\pi n}$,
т.е. $b_{2k-1}=\tfrac{4}{\pi(2k-1)}$. Значит
$\sign x=\tfrac4\pi\sum_{k\ge1}\tfrac{\sin(2k-1)x}{2k-1}$. В точке $x=\tfrac\pi2$ (гладкость, $f=1$)
$\sin\tfrac{(2k-1)\pi}2=(-1)^{k-1}$, откуда
\[
  1=\frac4\pi\sum_{k\ge1}\frac{(-1)^{k-1}}{2k-1}\ \Rightarrow\
  \frac\pi4=1-\frac13+\frac15-\frac17+\dots
\]

\prf{60}{Комплексная форма тригонометрического ряда Фурье}
\stmt В $L_2^{\C}[-\pi,\pi]$ система $\{e^{inx}\}$ ортогональна, $c_n=\tfrac1{2\pi}\int_{-\pi}^\pi
f e^{-inx}\dx$, и $c_n=\tfrac{a_n-ib_n}{2}$, $c_{-n}=\tfrac{a_n+ib_n}{2}$.

\dok Скалярное произведение $(f,g)=\int_{-\pi}^\pi f\bar g\dx$ эрмитово ($(f,g)=\overline{(g,f)}$),
поэтому $(f,f)=\int\abs f^2\ge0$ — сопряжение обеспечивает неотрицательность нормы.

\emph{Ортогональность:}
\[
  (e^{inx},e^{imx})=\int_{-\pi}^\pi e^{i(n-m)x}\dx=
  \begin{cases}0,&n\ne m,\\ 2\pi,&n=m,\end{cases}
\]
значит $\norm{e^{inx}}^2=2\pi$ и $c_n=\tfrac{(f,e^{inx})}{2\pi}=\tfrac1{2\pi}\int_{-\pi}^\pi
f(x)e^{-inx}\dx$.

\emph{Связь с вещественными:} подставив $e^{-inx}=\cos nx-i\sin nx$,
$c_n=\tfrac1{2\pi}\int f(\cos nx-i\sin nx)=\tfrac{a_n-ib_n}{2}$, и аналогично
$c_{-n}=\tfrac{a_n+ib_n}{2}$, $c_0=\tfrac{a_0}{2}$.

\end{document}
